\chapter{KURAMSAL ARKA PLAN}

\section{Hesaplama Grafikleri ve Operatör Füzyonu}
ONNX formatına aktarılan derin öğrenme modelleri, operatör düğümlerinden oluşan yönlü çevrimsiz grafikler (DAG) olarak ifade edilir.
Operatör füzyonu; ardışık operatörlerin tek bir birleşik çekirdek (kernel) veya sağlayıcıya özgü altgrafik olarak çalıştırılmasını hedefler.
Amaç ara tensörlerin belleğe yazılıp tekrar okunmasını azaltmak ve çalışma zamanı çağrı/dispatch ek yükünü düşürmektir.
Buna karşın füzyon, kernel seçimini değiştirebilir, bazı düşük seviye optimizasyonları engelleyebilir veya EP tarafında ek derleme maliyeti doğurabilir.

\section{ONNX Runtime Grafik Optimizasyon Seviyeleri}
ORT, oturum (session) oluşturulurken bir dizi grafik dönüşümü uygular.
Bu projede belgelenmiş ayarlar üzerinden “optimizasyon kapalı” konfigürasyonu ile “extended optimizasyon açık” konfigürasyonu karşılaştırılmıştır \cite{onnxruntime}.
Yüksek seviyede dönüşümler şu hedefleri taşır:
\begin{itemize}
	\item Gereksiz düğümleri kaldırma ve sabit ifadeleri katlama (constant folding).
	\item Kalıp (pattern) yeniden yazımlarıyla daha verimli grafik üretme.
	\item Füzyonları ve hedef çekirdeklere uygun düzen (layout) dönüşümlerini etkinleştirme.
\end{itemize}

\section{Execution Provider (EP) Kavramı ve OpenVINO EP}
ORT, yürütmeyi Execution Provider (EP) bileşenlerine devreder.
OpenVINO EP, Intel hedefleri için altgrafikleri derleyebilir; derleme başarısız olduğunda EP içi düşüş (fallback) veya CPU EP'ye düşüş gözlenebilir \cite{openvino_ep}.
Bu çalışmada NPU sınıfı hızlandırıcı davranışına yaklaşmak amacıyla öncelikli sağlayıcı olarak OpenVINO EP seçilmiştir.

\section{Profiling ve Operatör-Seviyesi Kanıt}
ORT profiling çıktısı, çekirdek bazlı süreleri ve metadata'yı içeren bir JSON izidir.
Bu iz üzerinde operatör türüne göre toplulaştırma yaparak baseline/optimized karşılaştırmasında hangi operatörlerin toplam süreyi domine ettiğini ölçmek mümkündür.
Bu ayrıştırma önemlidir; çünkü grafik yapısında sadeleşme olsa bile toplam gecikme değişmeyebilir.

\section{Memory Wall ve Roofline Sezgisi}
Pek çok modelde performans, aritmetik yoğunluktan ziyade veri taşımaya (bellek bant genişliğine) bağlıdır.
Memory wall \cite{memorywall} yaklaşımı, düğüm sayısındaki azalmanın neden her zaman hızlanmaya dönüşmediğini açıklamada kullanışlıdır.
Bu nedenle bulgular, basit bir “bellek-sınırlı / hesaplama-sınırlı” sezgisiyle yorumlanmıştır.